\subsection{How the Church Came to Keep Advent}\label{sec:history}

Christmas came before Advent. A Roman almanac of the year 354 already sets down the feast on the date the
Church still keeps: \latin{VIII kal. Ian. natus Christus in Betleem Iudeae},
``On the eighth day before the Kalends of January, Christ born in Bethlehem of
Judaea.'' Once the Church had a birthday of the Lord, she began to make ready for
it, and the making ready took different shapes in different lands.

In Spain a council of bishops met at Saragossa about the year 380. Its canons
survive, and one of them concerns the end of December. Its sense, in the
English of Hefele's history of the councils, is this:

\begin{quote}
(4) From the 17th December to the Feast of the Epiphany every one must attend
the church daily, and may not go with bare feet.\par
{\smallskip\noindent\small ---Council of Saragossa, canon 4, in the English of Hefele's history\par}
\end{quote}

\noindent The canon's span runs from 17 December to the Epiphany. It begins on
the date on which the last days of Advent now begin.

In Gaul the preparation was a fast. St Gregory of Tours, in his History of the
Franks, tells of Perpetuus, one of the bishops of Tours before him, who
``instituted the fasts and vigils to be observed during the course of the
year.'' Gregory copies out the list as his church still kept it, and among the
fasts is this: \latin{De depositione domni Martini usque natale Domini terna in
septimana ieiunia}, ``From the burial of our
lord Martin to the Nativity of our Lord, three fasts each week.'' St Martin was
buried on 11 November, so the fast of Tours ran some six weeks. A council at
Mâcon, which Hefele dates to 581, made the same rule:

\begin{quote}
\latin{Ut a feria sancti Martini usque natale Domini secunda, quarta et sexta
sabbati ieiunetur et sacrificia quadragensimali debeant ordine caelebrari.}\par
{\smallskip\noindent\small ---Council of Mâcon, canon 9\par}
\end{quote}

\noindent In the English of Hefele's history: ``From S. Martin's Day until Christmas, a fast
must be kept on Monday, Wednesday, and Friday of each week \ldots{} The
Sacrifice must be offered after the manner of the Lenten season.'' Dom
Prosper Guéranger, the abbot of Solesmes who wrote the great nineteenth-century
commentary on the liturgical year, counted these days, forty-three of them,
given to penance ``as though it were a second Lent, though less strict and severe
than that which precedes Easter.''

Rome kept a fast in December as well, but it was another thing. Pope St Leo the
Great preached on the four seasonal fasts of the Roman Church, which later ages
called the Ember days: in spring in Lent, in summer at Pentecost, in autumn in
the seventh month, and in winter in the tenth month, counting from March, which
is December. Leo's winter fast belonged to the turning of the seasons; it was
not an Advent fast.

Rome's preparation is found instead in her books of Mass prayers. The Gelasian
Sacramentary gives five Masses under the heading \latin{Orationes de Adventum
Domini}, prayers for the Lord's Advent, and the first of them opens with a cry: \latin{Excita, Domine, potentiam tuam, et
veni}, ``Stir up thy power, O Lord, and come.'' The Gregorian Sacramentary, in
the copy called the Hadrianum, has four Sundays. The last of them follows the
Ember Saturday, and the book marks it with three words, \latin{Die dominica
uacat}: the Sunday is vacant, without a station of its own. The station was the
church of Rome appointed for the day's Mass, and the same book names them, St
Mary Major for the Ember Wednesday, St Peter's for the Third Sunday. In the ninth century Amalarius of Metz, who wrote at length on the
meaning of the rites, found both reckonings in the books before him. The
author of the Missal, or of the Antiphonary, he says, urges us to make this
preparation through four weeks; but in old missals and lectionaries he found written
\latin{Hebdomada quinta ante Natalem Domini}, the fifth week before the
Nativity of the Lord. The church of Milan, whose rite bears the name of St
Ambrose, kept six Sundays of Advent. Rome's four
weeks prevailed in the West.

The same Amalarius had seen the Gloria fall silent in Advent:
\latin{Vidi tempore prisco, Gloria in excelsis Deo, prætermitti in diebus
Adventus Domini, et in aliquibus locis dalmaticas}. He had seen, in former
times, the angels' hymn left out in the days of Advent, and in some places the
dalmatic, the deacon's vestment, as well; and he thought there was reason in
it, for it called to mind the time of the Old Testament.

Later interpreters of the liturgy took up his reason and filled it out.
Honorius of Autun counted the four weeks as the ages of that waiting: the Gloria and the dalmatic are set aside
\latin{ad insinuandum tempus Veteris Testamenti, quod fluxit sub quatuor
dignitatibus Christi Nativitatem præfigurantibus, videlicet sub patriarchis, sub
judicibus, sub regibus, sub prophetis}, ``to suggest the time of the Old
Testament, which ran its course under four orders that prefigured the Nativity
of Christ, namely under the patriarchs, the judges, the kings and the
prophets.'' John Beleth, a twelfth-century master, noticed what Advent did
not give up.
The \latin{Te Deum laudamus}, ``We praise thee, O God,'' the old hymn of praise
that ends the night Office on festal days, is left out, he wrote, and the Gloria
too; but \latin{cum tempus istud partim sit lætitiæ, non subticetur Alleluia},
``since this season is partly one of gladness, the Alleluia is not silenced.''
And among his reasons for silencing the \latin{Te
Deum} he gave one any child could grasp: it is kept back \latin{quia nondum præsens est quem exspectamus},
``because he whom we await is not yet present.'' Advent was a season of joy held
in check long before the reform said so.

The Second Vatican Council asked for a revision of the liturgical year in which
the customs and observances handed down for the sacred seasons would be kept,
or brought back, as the conditions of our age allow, and each season would hold
its own native character and so feed the devotion of the faithful
(\emph{Sacrosanctum Concilium}, n.~107). The work was done within a few
years. Pope St Paul VI, by his letter \emph{Mysterii Paschalis} of 1969, approved
the Universal Norms on the Liturgical Year and the new General Roman Calendar,
to take effect on 1 January 1970. His apostolic constitution \emph{Missale Romanum} of 3 April 1969 set the
Sunday readings in a cycle of three years, so that the people might hear far
more of Scripture, and gave each weekday of Advent a prayer of its own; and it
ordered that the new Missal take effect on 30 November 1969, the First Sunday
of Advent.

Its prayers were gathered from the old books. The Gelasian and Gregorian
Sacramentaries gave many of them, and several more came from the Rotulus of
Ravenna, a roll of early Latin prayers for the days before Christmas. The
Lectionary for Mass appeared in 1969, and its second edition, with a fuller
introduction, in 1981. Paul VI wrote in 1974, in his apostolic exhortation
\emph{Marialis cultus}, that the Advent liturgy remembers the Mother of the
Lord most of all from 17 to 24 December and on the Sunday before Christmas, and
he endorsed the view that Advent is a fitting season for honouring her
(nn.~3--4).

The third edition of the Missal was published in 2002; a reprint with
emendations in 2008 added to the general calendar St Juan Diego on 9 December
and Our Lady of Guadalupe on 12 December. In the United States its English
translation came into use on the First Sunday of Advent, 27 November 2011. In
2019 the optional memorial of Our Lady of Loreto was added to the general
calendar on 10 December.
